\documentclass[12pt,a4paper]{article}
\usepackage[utf8]{inputenc}
\usepackage[T1]{fontenc}
\usepackage[french]{babel}
\usepackage{amsmath,amssymb,amsfonts}
\usepackage{multicol}
\usepackage{xcolor}
\usepackage{tikz}
\usepackage{tkz-tab}

% Configuration de la couleur rouge pour les titres et résultats
\definecolor{myred}{RGB}{180, 0, 0}
\newcommand{\redtitle}[1]{\textcolor{myred}{\textbf{#1}}}
\newcommand{\redbox}[1]{\color{myred}\boxed{#1}\color{black}}

\begin{document}

\section*{\underline{\redtitle{Correction de l'Exercice 1 : (5 pts)}}}

\subsection*{\underline{\redtitle{Partie A}}}

On considère le polynôme $P(x) = x^3 - 2x^2 - x + 2$.

\begin{enumerate}
    \item \textbf{Calculons $P(1)$ et concluons :}
    
    $P(1) = 1^3 - 2(1)^2 - 1 + 2 = 1 - 2 - 1 + 2 = 0$.
    
    \textbf{Conclusion :} Puisque $P(1) = 0$, alors $1$ est une racine du polynôme $P(x)$. Par conséquent, $P(x)$ est divisible par $(x - 1)$. \hfill \textbf{0,25pt}
    
    \item \textbf{Déterminons les réels $a$, $b$ et $c$ :}
    
    Développons l'expression $(x - 1)(ax^2 + bx + c)$ :
    \begin{align*}
        (x - 1)(ax^2 + bx + c) &= ax^3 + bx^2 + cx - ax^2 - bx - c \\
        &= ax^3 + (b - a)x^2 + (c - b)x - c
    \end{align*}
    Par identification avec $P(x) = x^3 - 2x^2 - x + 2$, on obtient le système suivant :
    \[
    \begin{cases}
        a = 1 \\
        b - a = -2 \\
        c - b = -1 \\
        -c = 2
    \end{cases}
    \implies
    \begin{cases}
        a = 1 \\
        b = -2 + 1 = -1 \\
        c = -2
    \end{cases}
    \]
    Les réels cherchés sont donc : $\redbox{a = 1 \,;~ b = -1 \,;~ c = -2}$.
    
    Ainsi, $P(x) = (x - 1)(x^2 - x - 2)$. \hfill \textbf{0,5pt}
    
    \item On pose $Q(x) = x^2 - x - 2$.
    \begin{enumerate}
        \item \textbf{Donner une factorisation complète de $P(x)$ :}
        
        Cherchons les racines du trinôme $Q(x) = x^2 - x - 2$. \\
        On remarque une racine évidente $x_1 = -1$ (car $(-1)^2 - (-1) - 2 = 1 + 1 - 2 = 0$). \\
        Le produit des racines étant $\frac{c}{a} = -2$, la seconde racine est $x_2 = 2$.
        
        D'où la factorisation de $Q(x) = (x + 1)(x - 2)$.
        
        La factorisation complète de $P(x)$ est donc :
        \[\redbox{P(x) = (x - 1)(x + 1)(x - 2)}\] \hfill \textbf{0,25pt}
        
        \item \textbf{Résolvons dans $\mathbb{R}$ l'équation $P(x) = 0$ et l'inéquation $P(x) \leq 0$ :}
        
        \begin{itemize}
            \item \textbf{Équation $P(x) = 0$ :}
            \[ P(x) = 0 \iff (x - 1)(x + 1)(x - 2) = 0 \]
            L'ensemble des solutions est : $\redbox{S = \{-1 \,;~ 1 \,;~ 2\}}$.
            
            \item \textbf{Inéquation $P(x) \leq 0$ :} \\
            Dressons le tableau de signes de $P(x)$ à l'aide de ses facteurs :
            
            \begin{center}
            \begin{tikzpicture}
                \tkzTabInit[lgt=3, espcl=1.5]{$x$ / 1 , $x+1$ / 0.6, $x-1$ / 0.6, $x-2$ / 0.6, $P(x)$ / 0.8}{$-\infty$, $-1$, $1$, $2$, $+\infty$}
                \tkzTabLine{, -, z, +, t, +, t, +, }
                \tkzTabLine{, -, t, -, z, +, t, +, }
                \tkzTabLine{, -, t, -, t, -, z, +, }
                \tkzTabLine{, -, z, +, z, -, z, +, }
            \end{tikzpicture}
            \end{center}
            
            L'ensemble des solutions de l'inéquation $P(x) \le 0$ est : $\redbox{S = ]-\infty \,;~ -1] \cup [1 \,;~ 2]}$. \hfill \textbf{1pt}
        \end{itemize}
        
        \item \textbf{En déduire les solutions dans $\mathbb{R}$ de l'équation $P(x-2) = 0$ et de l'inéquation $P(x-2) \leq 0$ :}
        
        En posant $X = x - 2$, on utilise directement les résultats de la question précédente :
        
        \begin{itemize}
            \item \textbf{Pour l'équation $P(x-2) = 0$ :}
            \[
            x - 2 = -1 \implies x = 1
            \]
            \[
            x - 2 = 1 \implies x = 3
            \]
            \[
            x - 2 = 2 \implies x = 4
            \]
            L'ensemble des solutions est : $\redbox{S_1 = \{1 \,;~ 3 \,;~ 4\}}$.
            
            \item \textbf{Pour l'inéquation $P(x-2) \leq 0$ :}
            \[ X \in ]-\infty \,;~ -1] \cup [1 \,;~ 2] \]
            Ce qui se traduit par :
            \begin{align*}
                x - 2 \le -1 &\implies x \le 1 \\
                \text{ou } 1 \le x - 2 \le 2 &\implies 3 \le x \le 4
            \end{align*}
            L'ensemble des solutions est : $\redbox{S_2 = ]-\infty \,;~ 1] \cup [3 \,;~ 4]}$. \hfill \textbf{1pt}
        \end{itemize}
    \end{enumerate}
\end{enumerate}

\subsection*{\underline{\redtitle{Partie B}}}

Soit l'équation $(E) : x^2 - 2x - 3 = 0$ de la forme $ax^2 + bx + c = 0$ avec $a=1$, $b=-2$ et $c=-3$.

\begin{enumerate}
    \item \redtitle{Calculons sans calculer $x'$ et $x''$ :}
    
    On sait que la somme des racines est $S = -\frac{b}{a}$ et le produit est $P = \frac{c}{a}$.
    
    \begin{multicols}{2}
    \setlength{\columnseprule}{0.1mm}
    \begin{enumerate}
        \item $A = x' + x'' = -\dfrac{-2}{1}$ \\
        Donc : $\redbox{A = 2}$ \hfill \textbf{0,25pt}
        \item $B = x'x'' = \dfrac{-3}{1}$ \\
        Donc : $\redbox{B = -3}$ \hfill \textbf{0,25pt}
    \end{enumerate}
    \end{multicols}

    \item \redtitle{Déduisons-en les valeurs de :}
    
    \begin{enumerate}
        \item Réduisons au même dénominateur :
        \[ A' = \dfrac{1}{x'} + \dfrac{1}{x''} = \dfrac{x'' + x'}{x'x''} = \dfrac{A}{B} \]
        \[ \redbox{A' = -\dfrac{2}{3}} \] \hfill \textbf{0,5pt}
        
        \item Identité remarquable :
        \begin{align*}
            C &= {x'}^2 + {x''}^2 \\
              &= (x' + x'')^2 - 2x'x'' \\
              &= A^2 - 2B \\
              &= (2)^2 - 2(-3) \\
              &= 4 + 6
        \end{align*}
        \[ \redbox{C = 10} \] \hfill \textbf{0,5pt}
        
        \item Identité remarquable :
        \begin{align*}
            D &= {x'}^3 + {x''}^3 \\
              &= (x' + x'')( {x'}^2 - x'x'' + {x''}^2 ) \\
              &= A(C - B) \\
              &= 2(10 - (-3)) \\
              &= 2(13)
        \end{align*}
        \[ \redbox{D = 26} \] \hfill \textbf{0,5pt}
    \end{enumerate}

\end{enumerate}


\section*{\underline{\redtitle{Correction de l'Exercice 2 : (3 pts)}}}

Soit l'équation $(E) : x^2 - (2m + 1)x + m^2 - 2 = 0$ de la forme $ax^2 + bx + c = 0$ avec :
$a = 1$, $b = -(2m + 1)$ et $c = m^2 - 2$.

\begin{enumerate}
    \item \redtitle{Discutons suivant les valeurs de $m$, les solutions de l'équation $(E)$ :}
    
    Calculons le discriminant $\Delta$ de l'équation $(E)$ :
    \begin{align*}
        \Delta &= b^2 - 4ac \\
               &= [-(2m + 1)]^2 - 4(1)(m^2 - 2) \\
               &= (4m^2 + 4m + 1) - 4m^2 + 8 \\
               &= 4m + 9
    \end{align*}
    
    Étudions le signe de $\Delta$ :
    \[ \Delta = 0 \iff 4m + 9 = 0 \iff m = -\dfrac{9}{4} \]
    
    \textbf{Discussion :}
    \begin{itemize}
        \item Si $m \in \left]-\infty \,;~ -\dfrac{9}{4}\right[$, alors $\Delta < 0$ : L'équation $(E)$ n'admet aucune solution réelle. $\redbox{S = \emptyset}$.
        \item Si $m = -\dfrac{9}{4}$, alors $\Delta = 0$ : L'équation $(E)$ admet une solution double :
        \[ x_0 = -\dfrac{b}{2a} = \dfrac{2\left(-\frac{9}{4}\right) + 1}{2} = \dfrac{-\frac{9}{2} + 1}{2} = -\dfrac{7}{4} \]
        L'ensemble des solutions est : $\redbox{S = \left\{-\dfrac{7}{4}\right\}}$.
        \item Si $m \in \left]-\dfrac{9}{4} \,;~ +\infty\right[$, alors $\Delta > 0$ : L'équation $(E)$ admet deux solutions distinctes :
        \[ \redbox{x' = \dfrac{(2m+1) - \sqrt{4m+9}}{2} \quad \text{et} \quad x'' = \dfrac{(2m+1) + \sqrt{4m+9}}{2}} \]
    \end{itemize} \hfill \textbf{(0,75 pt)}
    
    \item \redtitle{Trouvons une relation indépendante de $m$ entre $x'$ et $x''$ :}
    
    Exprimons la somme $S$ et le produit $P$ des racines en fonction de $m$ :
    \[ S = x' + x'' = -\dfrac{b}{a} = 2m + 1 \implies m = \dfrac{S - 1}{2} \]
    \[ P = x'x'' = \dfrac{c}{a} = m^2 - 2 \]
    
    En substituant l'expression de $m$ dans celle de $P$, on obtient :
    \[ P = \left(\dfrac{S - 1}{2}\right)^2 - 2 \iff P = \dfrac{S^2 - 2S + 1}{4} - 2 \]
    \[ 4P = S^2 - 2S + 1 - 8 \iff S^2 - 2S - 4P - 7 = 0 \]
    
    En remplaçant par $x'$ et $x''$, la relation indépendante de $m$ est :
    \[ \redbox{(x' + x'')^2 - 2(x' + x'') - 4x'x'' - 7 = 0} \] \hfill \textbf{(0,5 pt)}
    
    \item Pour que $(E)$ admette deux solutions, il faut d'abord que celles-ci existent, c'est-à-dire $\Delta \geq 0 \iff m \geq -\dfrac{9}{4}$.

la somme $S = 2m+1$ et du produit $P = m^2-2$.
    
    Dressons un tableau récapitulatif des signes pour mener la discussion :
    
    \begin{center}
    \begin{tikzpicture}
        \tkzTabInit[lgt=3, espcl=1.5]{$m$ / 1 , $\Delta=4m+9$ / 0.6, $P=m^2-2$ / 0.6, $S=2m+1$ / 0.6}{$-\dfrac{9}{4}$, $-\sqrt{2}$, $-\dfrac{1}{2}$, $\sqrt{2}$, $+\infty$}
        \tkzTabLine{z, +, t, +, t, +, t, +, }
        \tkzTabLine{, +, z, -, t, -, z, +, }
        \tkzTabLine{, -, t, -, z, +, t, +, }
    \end{tikzpicture}
    \end{center}
    
    \begin{enumerate}
        \item \redtitle{Deux solutions de signes contraires :}
        
        L'équation admet deux solutions de signes contraires si et seulement si leur produit est strictement négatif ($P < 0$). 
        
        En lisant sur la ligne de $m^2-2$ dans notre tableau sur l'intervalle d'existence ($\Delta \geq 0$), on obtient :
        \[ \redbox{m \in ]-\sqrt{2} \,;~ \sqrt{2}[} \] \hfill \textbf{(0,25 pt)}
        
        \item \redtitle{Deux solutions positives :}
        
        L'équation admet deux solutions positives si et seulement si :
        \[ \begin{cases} \Delta \geq 0 \\ P > 0 \\ S > 0 \end{cases} \]
        D'après le tableau, les trois expressions ($4m+9$, $m^2-2$ et $2m+1$) sont simultanément strictement positives sur le dernier intervalle.
        
        La condition est donc : $\redbox{m \in ]\sqrt{2} \,;~ +\infty[}$. \hfill \textbf{(0,75 pt)}
        
        \item \redtitle{Deux solutions négatives :}
        
        L'équation admet deux solutions négatives si et seulement si :
        \[ \begin{cases} \Delta \geq 0 \\ P > 0 \\ S < 0 \end{cases} \]
        D'après le tableau, sur le premier intervalle d'existence, $\Delta \geq 0$ (le signe de $4m+9$ est positif ou nul), $P > 0$ (le signe de $m^2-2$ est positif) et $S < 0$ (le signe de $2m+1$ est négatif).
        
        La condition est donc : $\redbox{m \in \left[ -\dfrac{9}{4} \,;~ -\sqrt{2} \right[}$. \hfill \textbf{(0,75 pt)}
    \end{enumerate}
\end{enumerate}


\section*{\underline{\redtitle{Correction de l'Exercice 3 : (6 pts)}}}

\begin{enumerate}
    \item \redtitle{Résolvons dans $\mathbb{R}$ les équations suivantes :}
    
    \begin{enumerate}
        \item $3x^2 - 5x + 2 = 0$
        
        On remarque que $a + b + c = 3 - 5 + 2 = 0$. 
        
        Les racines sont donc évidentes : $x_1 = 1$ et $x_2 = \dfrac{c}{a} = \dfrac{2}{3}$.
        
        L'ensemble des solutions est : $\redbox{S = \left\{\dfrac{2}{3} \,;~ 1\right\}}$. \hfill \textbf{(0,25 pt)}
        
        \item $2|x - 2|^2 + 7|x - 2| + 3 = 0$
        
        Posons $X = |x - 2|$ avec $X \geq 0$.
        
        L'équation devient :$2X^2 + 7X + 3 = 0$.
        
        Calculons le discriminant :
        
         $\Delta = 7^2 - 4(2)(3) = 49 - 24 = 25 = 5^2$.
        
        \[ X_1 = \dfrac{-7 - 5}{4} = -3 \quad \text{et} \quad X_2 = \dfrac{-7 + 5}{4} = -\dfrac{2}{4} = -\dfrac{1}{2} \]
        
        Comme $X = |x - 2|$ doit être positif ou nul ($X \geq 0$), les deux valeurs $X_1$ et $X_2$ sont rejetées.
        
        L'ensemble des solutions est : $\redbox{S = \emptyset}$. \hfill \textbf{(0,75 pt)}
        
        \item $\left(x + \dfrac{1}{x}\right)^2 = 4\left(x + \dfrac{1}{x}\right) - 3$ \quad (Condition d'existence : $x \neq 0$)
        
        Posons $X = x + \dfrac{1}{x}$. L'équation devient :
        \[ X^2 = 4X - 3 \iff X^2 - 4X + 3 = 0 \]
        Ici, $1 - 4 + 3 = 0$, donc les solutions en $X$ sont $X_1 = 1$ et $X_2 = 3$.
        
        Revenons à $x$ :
        \begin{itemize}
            \item \textbf{Si $X = 1$ :} $x + \dfrac{1}{x} = 1 \iff x^2 - x + 1 = 0$. 
            
            $\Delta = (-1)^2 - 4(1)(1) = -3 < 0 \implies$ Pas de solution réelle.
            \item \textbf{Si $X = 3$ :} $x + \dfrac{1}{x} = 3 \iff x^2 - 3x + 1 = 0$.
            
            $\Delta = (-3)^2 - 4(1)(1) = 9 - 4 = 5$.
            \[ x_1 = \dfrac{3 - \sqrt{5}}{2} \quad \text{et} \quad x_2 = \dfrac{3 + \sqrt{5}}{2} \]
        \end{itemize}
        L'ensemble des solutions est : 
        
        $\redbox{S = \left\{\dfrac{3 - \sqrt{5}}{2} \,;~ \dfrac{3 + \sqrt{5}}{2}\right\}}$. \hfill \textbf{(0,75 pt)}
    \end{enumerate}

    \item \redtitle{Résolvons dans $\mathbb{R}$ les inéquations suivantes :}
    
    \begin{enumerate}
        \item $x^2 - x + 6 \geq 0$
        
        Calculons le discriminant : 
        
        $\Delta = (-1)^2 - 4(1)(6) = 1 - 24 = -23 < 0$.
        
        Le trinôme n'a pas de racine réelle, il est donc strictement du signe de $a=1$ (positif) pour tout réel $x$.
        
        L'ensemble des solutions est : $\redbox{S = \mathbb{R}}$. \hfill \textbf{(0,25 pt)}
        
        \item $\dfrac{2x^2 + 3x - 5}{-x^2 + 3x + 10} \leq 0$
        
        Trouvons les racines du numérateur $N(x) = 2x^2 + 3x - 5$ :
        $a+b+c = 2+3-5=0 \implies x_1 = 1$ et $x_2 = -\dfrac{5}{2}$.
        
        Trouvons les racines du dénominateur $D(x) = -x^2 + 3x + 10$ :
        $\Delta = 3^2 - 4(-1)(10) = 9 + 40 = 49 = 7^2$.
        \[ x_3 = \dfrac{-3 - 7}{-2} = 5 \quad \text{et} \quad x_4 = \dfrac{-3 + 7}{-2} = -2 \]
        Les valeurs interdites sont donc $-2$ et $5$. Dressons le tableau de signes :
        
        \begin{center}
        \begin{tikzpicture}
            \tkzTabInit[lgt=3, espcl=1.2]{$x$ / 1 , $N(x)$ / 0.6, $D(x)$ / 0.6, Quotient / 0.8}{$-\infty$, $-\dfrac{5}{2}$, $-2$, $1$, $5$, $+\infty$}
            \tkzTabLine{, +, z, -, t, -, z, +, t, +, }
            \tkzTabLine{, -, t, -, z, +, t, +, z, -, }
            \tkzTabLine{, -, z, +, d, -, z, +, d, -, }
        \end{tikzpicture}
        \end{center}
        
        L'ensemble des solutions est : 
        
        $\redbox{S = \left[ -\infty \,;~ -\dfrac{5}{2} \right] \cup ]-2 \,;~ 1] \cup ]5 \,;~ +\infty[}$. \hfill \textbf{(0,75 pt)}
    \end{enumerate}

    \item \redtitle{Résolvons dans $\mathbb{R}^2$ les systèmes suivants :}
    
    \begin{enumerate}
        \item $\begin{cases} x + y = 3 \\ x^2 + 3xy + y^2 = 11 \end{cases}$
        
        On sait que $x^2 + y^2 = (x+y)^2 - 2xy$. La deuxième équation devient :
        \[ (x+y)^2 - 2xy + 3xy = 11 \iff (x+y)^2 + xy = 11 \]
        En substituant $x+y=3$, on obtient :
        
         $3^2 + xy = 11 \iff 9 + xy = 11 \iff xy = 2$.
        
        On se ramène à la recherche de deux nombres connaissant leur somme $S=3$ et leur produit $P=2$, qui sont solutions de 
        
        $t^2 - 3t + 2 = 0 \implies t_1 = 1$ et $t_2 = 2$.
        
        L'ensemble des solutions est :
        
         $\redbox{S = \{(1 \,;~ 2) \,;~ (2 \,;~ 1)\}}$. \hfill \textbf{(0,75 pt)}
        
        \item $\begin{cases} x^2 + y^2 = 5 \\ xy = 2 \end{cases}$
        
        On utilise les identités remarquables :
        \begin{itemize}
            \item $(x+y)^2 = x^2 + y^2 + 2xy = 5 + 2(2) = 9 \implies x + y = 3 \text{ ou } x + y = -3$
            \item $(x-y)^2 = x^2 + y^2 - 2xy = 5 - 2(2) = 1 \implies x - y = 1 \text{ ou } x - y = -1$
        \end{itemize}
        En combinant les cas ou en résolvant par somme/produit ($S=3, P=2$ ou $S=-3, P=2$), on trouve :
        \[ \redbox{S = \{(2 \,;~ 1) \,;~ (1 \,;~ 2) \,;~ (-2 \,;~ -1) \,;~ (-1 \,;~ -2)\}} \] \hfill \textbf{(0,75 pt)}
        
        \item $\begin{cases} x^3 + y^3 = 7 \\ x + y = 1 \end{cases}$
        
        On sait que $x^3 + y^3 = (x+y)(x^2 - xy + y^2) = (x+y)[(x+y)^2 - 3xy]$.
        La première équation devient :
        \[ 1 \cdot (1^2 - 3xy) = 7 \iff 1 - 3xy = 7 \iff -3xy = 6 \iff xy = -2 \]
        On cherche deux nombres de somme $S=1$ et produit $P=-2$, solutions de $t^2 - t - 2 = 0$.
        Les racines sont $t_1 = -1$ et $t_2 = 2$.
        
        L'ensemble des solutions est : 
        
        $\redbox{S = \{(-1 \,;~ 2) \,;~ (2 \,;~ -1)\}}$. \hfill \textbf{(0,75 pt)}
    \end{enumerate}

    \item \redtitle{En utilisant la méthode de Cramer, résolvons le système :} $\begin{cases} 3x + 4y = 5 \\ 2x - 3y = 1 \end{cases}$
    
    Calculons le déterminant principal $\det(S)$ :
    \[ \det(S) = \begin{vmatrix} 3 & 4 \\ 2 & -3 \end{vmatrix} = 3(-3) - 4(2) = -9 - 8 = -17 \]
    Comme $\det(S) \neq 0$, le système admet un couple solution unique.
    
    Calculons les déterminants secondaires $\det(x)$ et $\det(y)$ :
    \[ \det(x) = \begin{vmatrix} 5 & 4 \\ 1 & -3 \end{vmatrix} = 5(-3) - 4(1) = -15 - 4 = -19 \]
    \[ \det(y) = \begin{vmatrix} 3 & 5 \\ 2 & 1 \end{vmatrix} = 3(1) - 5(2) = 3 - 10 = -7 \]
    
    On en déduit les valeurs de $x$ et $y$ :
    \[ x = \dfrac{\det(x)}{\det(S)} = \dfrac{-19}{-17} = \dfrac{19}{17} \quad \text{et} \quad y = \dfrac{\det(y)}{\det(S)} = \dfrac{-7}{-17} = \dfrac{7}{17} \]
    
    L'ensemble des solutions est : $\redbox{S = \left\{\left(\dfrac{19}{17} \,;~ \dfrac{7}{17}\right)\right\}}$. \hfill \textbf{(0,75 pt)}
\end{enumerate}


\section*{\underline{\redtitle{Correction de l'Exercice 4 : (6 pts)}}}

Soit $ABC$ un triangle quelconque et $I$ le milieu de $[BC]$.

\begin{enumerate}
    \item \redtitle{Déterminons l'ensemble des valeurs réelles de $\alpha$ pour que le barycentre $G_\alpha$ existe :}
    
    Le système admet un unique barycentre si et seulement si la somme des coefficients est non nulle.
    \begin{align*}
        \sum \text{Coefficients} \neq 0 &\iff (2\alpha^2 - 3) + \alpha + (-\alpha) \neq 0 \\
        & \iff 2\alpha^2 - 3 \neq 0 \\
        & \iff \alpha^2 \neq \dfrac{3}{2} \\
        & \iff \alpha \neq -\sqrt{\dfrac{3}{2}} \quad \text{et} \quad \alpha \neq \sqrt{\dfrac{3}{2}}
    \end{align*}
    L'ensemble des valeurs de $\alpha$ recherché est : $\redbox{\mathbb{R} \setminus \left\{ -\sqrt{\dfrac{3}{2}} \,;~ \sqrt{\dfrac{3}{2}} \right\}}$. \hfill \textbf{(1pt)}
    
    \item \redtitle{Montrons que pour tout réel $\alpha$ autorisé, on a $\overrightarrow{AG_\alpha} = -\dfrac{\alpha}{2\alpha^2 - 3}\overrightarrow{BC}$ :}
    
    Par définition du barycentre $G_\alpha$, on a :
    \[ (2\alpha^2 - 3)\overrightarrow{G_\alpha A} + \alpha\overrightarrow{G_\alpha B} - \alpha\overrightarrow{G_\alpha C} = \overrightarrow{0} \]
    En introduisant le point $A$ à l'aide de la relation de Chasles dans les deux derniers vecteurs :
    \[ (2\alpha^2 - 3)\overrightarrow{G_\alpha A} + \alpha\left(\overrightarrow{G_\alpha A} + \overrightarrow{AB}\right) - \alpha\left(\overrightarrow{G_\alpha A} + \overrightarrow{AC}\right) = \overrightarrow{0} \]
    \[ (2\alpha^2 - 3)\overrightarrow{G_\alpha A} + \alpha\overrightarrow{G_\alpha A} - \alpha\overrightarrow{G_\alpha A} + \alpha\overrightarrow{AB} - \alpha\overrightarrow{AC} = \overrightarrow{0} \]
    \[ (2\alpha^2 - 3)\overrightarrow{G_\alpha A} + \alpha\left(\overrightarrow{AB} - \overrightarrow{AC}\right) = \overrightarrow{0} \]
    Sachant que $\overrightarrow{G_\alpha A} = -\overrightarrow{AG_\alpha}$ et $\overrightarrow{AB} - \overrightarrow{AC} = \overrightarrow{CB} = -\overrightarrow{BC}$, on obtient :
    \[ -(2\alpha^2 - 3)\overrightarrow{AG_\alpha} - \alpha\overrightarrow{BC} = \overrightarrow{0} \iff (2\alpha^2 - 3)\overrightarrow{AG_\alpha} = -\alpha\overrightarrow{BC} \]
    D'où la relation finale : $\redbox{\overrightarrow{AG_\alpha} = -\dfrac{\alpha}{2\alpha^2 - 3}\overrightarrow{BC}}$. \hfill \textbf{(1pt)}

    \item \redtitle{Construisons les points $G_1$ et $G_{-1}$ :}
    
    En utilisant la relation établie à la question (2) :
    \begin{itemize}
        \item Pour $\alpha = 1$ : $\overrightarrow{AG_1} = -\dfrac{1}{2(1)^2 - 3}\overrightarrow{BC} = -\dfrac{1}{-1}\overrightarrow{BC} \implies \redbox{\overrightarrow{AG_1} = \overrightarrow{BC}}$.
        \item Pour $\alpha = -1$ : $\overrightarrow{AG_{-1}} = -\dfrac{-1}{2(-1)^2 - 3}\overrightarrow{BC} = \dfrac{1}{-1}\overrightarrow{BC} \implies \redbox{\overrightarrow{AG_{-1}} = -\overrightarrow{BC}}$.
    \end{itemize}
    \textbf{Utilisation du point $I$ :} \\
    Puisque $I$ est le milieu de $[BC]$, on sait par ailleurs que
    
     $\overrightarrow{BC} = 2\overrightarrow{IC} = 2\overrightarrow{BI}$. On peut donc l'utiliser pour la précision géométrique si on le souhaite, le point $G_1$ complétant le parallélogramme $ABCG_1$. \hfill \textbf{(1pt)}

\begin{center}
\begin{tikzpicture}[scale=1.2, >=stealth]
    % Définition des points du triangle ABC
    \coordinate (A) at (2,3);
    \coordinate (B) at (0,0);
    \coordinate (C) at (5,0);
    
    % Calcul du milieu I de [BC]
    \coordinate (I) at (2.5,0);
    
    % Calcul des points G_1 et G_{-1}
    % G_1 tel que AG_1 = BC -> G_1 = A + (C - B) = (2+5, 3+0) = (7,3)
    \coordinate (G1) at (7,3);
    % G_{-1} tel que AG_{-1} = -BC -> G_{-1} = A - (C - B) = (2-5, 3-0) = (-3,3)
    \coordinate (G-1) at (-3,3);

    % Tracé du triangle ABC
    \draw[thick] (A) -- (B) -- (C) -- cycle;
    
    % Tracé des segments de construction vers G_1 et G_{-1}
    \draw[dashed, blue] (A) -- (G1);
    \draw[dashed, blue] (C) -- (G1);
    \draw[dashed, purple] (A) -- (G-1);
    \draw[dashed, purple] (B) -- (G-1);
    
    % Tracé de la ligne droite reliant G_{-1}, A et G_1
    \draw[thick, red] (G-1) -- (G1);

    % Marquage du milieu I (un petit trait vertical ou une croix)
    \draw[thick] (2.5,-0.1) -- (2.5,0.1);

    % Affichage des points (noeuds)
    \fill (A) circle (1.5pt) node[above] {$A$};
    \fill (B) circle (1.5pt) node[below left] {$B$};
    \fill (C) circle (1.5pt) node[below right] {$C$};
    \fill (I) circle (1.5pt) node[below] {$I$};
    \fill[red] (G1) circle (2pt) node[above right] {$G_1$};
    \fill[red] (G-1) circle (2pt) node[above left] {$G_{-1}$};

    % Optionnel : Vecteurs pour bien comprendre le déplacement
    \draw[->, thick, cyan] (B) -- (C) node[midway, below] {$\overrightarrow{BC}$};
    \draw[->, thick, cyan] (A) -- (G1) node[midway, above] {$\overrightarrow{BC}$};
    \draw[->, thick, orange] (A) -- (G-1) node[midway, above] {$-\overrightarrow{BC}$};

\end{tikzpicture}
\end{center}

    \item \redtitle{Montrons que $I$ est le milieu de $[G_1 G_{-1}]$ :}
    
    Calculons le vecteur $\overrightarrow{IG_1} + \overrightarrow{IG_{-1}}$ en passant par le point $A$ :
    \begin{align*}
        \overrightarrow{IG_1} + \overrightarrow{IG_{-1}} &= \left(\overrightarrow{IA} + \overrightarrow{AG_1}\right) + \left(\overrightarrow{IA} + \overrightarrow{AG_{-1}}\right) \\
        &= 2\overrightarrow{IA} + \overrightarrow{AG_1} + \overrightarrow{AG_{-1}}
    \end{align*}
    En remplaçant par les valeurs géométriques trouvées à la question (3) :
    \[ \overrightarrow{IG_1} + \overrightarrow{IG_{-1}} = 2\overrightarrow{IA} + \overrightarrow{BC} + (-\overrightarrow{BC}) = 2\overrightarrow{IA} \]
    \textit{Note géométrique :} Pour que $I$ soit le milieu de $[G_1G_{-1}]$, la somme vectorielle précédente doit valoir $\overrightarrow{0}$, ce qui n'est vrai que si $A=I$. Reprenons par une autre décomposition plus directe en partant de $G_1$ et $G_{-1}$ :
    
    On a d'une part : $\overrightarrow{AG_1} = \overrightarrow{BC} \implies \overrightarrow{G_1C} = \overrightarrow{AB}$ ($ABCG_1$ est un parallélogramme). \\
    D'autre part : $\overrightarrow{AG_{-1}} = \overrightarrow{CB} \implies \overrightarrow{G_{-1}B} = \overrightarrow{AC}$.
    
    Calculons le milieu du segment $[G_1G_{-1}]$. 
    
    On sait que $\overrightarrow{AG_1} + \overrightarrow{AG_{-1}} = \overrightarrow{BC} - \overrightarrow{BC} = \overrightarrow{0}$. \\
    Cela signifie que $A$ est en réalité le milieu de $[G_1 G_{-1}]$.
    
    Vérifions si l'énoncé d'origine ne contenait pas une coquille textuelle entre le point $A$ et le point $I$. Puisque $\overrightarrow{AG_1} + \overrightarrow{AG_{-1}} = \overrightarrow{0}$, on en déduit formellement que :
    \[ \redbox{A \text{ est le milieu de } [G_1G_{-1}]} \] \hfill \textbf{(1pt)}

    \item \redtitle{Déterminons l'ensemble (E) des points M du plan dans chaque cas :}
    
    Pour simplifier le vecteur commun aux deux membres
    
     $-\overrightarrow{MA} + \overrightarrow{MB} - \overrightarrow{MC}$, remarquons que la somme des coefficients est $-1 + 1 - 1 = -1 \neq 0$. Introduisons le barycentre particulier de ce système, qui correspond à $G_\alpha$ pour $\alpha = 1$,
     
      c'est-à-dire $G_1 = \{(A, -1); (B, 1); (C, -1)\}$. 
    On a pour tout point $M$ :
    \[ -\overrightarrow{MA} + \overrightarrow{MB} - \overrightarrow{MC} = -\overrightarrow{MG_1} \]
    
    \begin{enumerate}
        \item \textbf{Cas 1 :} $\left\|-\overrightarrow{MG_1}\right\| = 9$
        \[ \iff MG_1 = 9 \]
        L'ensemble $(E_a)$ est le \textbf{cercle} de centre $G_1$ et de rayon $r = 9$. \hfill \textbf{(1pt)}
        
        \item \textbf{Cas 2 :} $\left\|-\overrightarrow{MG_1}\right\| = \left\|\overrightarrow{MA} - \overrightarrow{MB}\right\|$
        
        Le second membre correspond à un vecteur constant indépendant de $M$ : 
        $\overrightarrow{MA} - \overrightarrow{MB} = \overrightarrow{BA}$.
        L'équation devient donc :
        \[ MG_1 = BA \]
        L'ensemble $(E_b)$ est le \textbf{cercle} de centre $G_1$ et de rayon $r = BA$. \hfill \textbf{(1pt)}
    \end{enumerate}

\end{enumerate}

\end{document}